\documentclass[11pt,twoside,final,openright]{report}
\usepackage{a4,graphicx,html,parskip,setspace,times,xspace}
\setstretch{1.15}


\def\Xend{{Xend}\xspace}
\def\xend{{xend}\xspace}

\latexhtml{\newcommand{\path}[1]{{\small {\tt #1}}}}{\newcommand{\path}[1]{{\tt #1}}}



\begin{document}

% TITLE PAGE
\pagestyle{empty}
\begin{center}
\vspace*{\fill}
\includegraphics{figs/xenlogo.eps}
\vfill
\vfill
\vfill
\begin{tabular}{l}
{\Huge \bf Users' manual} \\[4mm]
{\huge Xen v2.0 for x86} \\[80mm]

{\Large Xen is Copyright (c) 2002-2004, The Xen Team} \\[3mm]
{\Large University of Cambridge, UK} \\[20mm]
\end{tabular}
\end{center}

{\bf
DISCLAIMER: This documentation is currently under active development
and as such there may be mistakes and omissions --- watch out for
these and please report any you find to the developer's mailing list.
Contributions of material, suggestions and corrections are welcome.
}

\vfill
\cleardoublepage

% TABLE OF CONTENTS
\pagestyle{plain}
\pagenumbering{roman}
{ \parskip 0pt plus 1pt
  \tableofcontents }
\cleardoublepage

% PREPARE FOR MAIN TEXT
\pagenumbering{arabic}
\raggedbottom
\widowpenalty=10000
\clubpenalty=10000
\parindent=0pt
\parskip=5pt
\renewcommand{\topfraction}{.8}
\renewcommand{\bottomfraction}{.8}
\renewcommand{\textfraction}{.2}
\renewcommand{\floatpagefraction}{.8}
\setstretch{1.1}

\part{Introduction and Tutorial}
\chapter{Introduction}

Xen is a {\em paravirtualising} virtual machine monitor (VMM), or
`hypervisor', for the x86 processor architecture.  Xen can securely
execute multiple virtual machines on a single physical system with
close-to-native performance.  The virtual machine technology
facilitates enterprise-grade functionality, including:

\begin{itemize}
\item Virtual machines with performance close to native
  hardware.
\item Live migration of running virtual machines between physical hosts.
\item Excellent hardware support (supports most Linux device drivers).
\item Sandboxed, restartable device drivers.
\end{itemize}

Paravirtualisation permits very high performance virtualisation,
even on architectures like x86 that are traditionally
very hard to virtualise.
The drawback of this approach is that it requires operating systems to
be {\em ported} to run on Xen.  Porting an OS to run on Xen is similar
to supporting a new hardware platform, however the process
is simplified because the paravirtual machine architecture is very
similar to the underlying native hardware. Even though operating system
kernels must explicitly support Xen, a key feature is that user space
applications and libraries {\em do not} require modification.

Xen support is available for increasingly many operating systems:
right now, Linux 2.4, Linux 2.6 and NetBSD are available for Xen 2.0.
A FreeBSD port is undergoing testing and will be incorporated into the
release soon. Other OS ports, including Plan 9, are in progress.  We
hope that that arch-xen patches will be incorporated into the
mainstream releases of these operating systems in due course (as has
already happened for NetBSD).

Possible usage scenarios for Xen include:
\begin{description}
\item [Kernel development.] Test and debug kernel modifications in a
      sandboxed virtual machine --- no need for a separate test
      machine.
\item [Multiple OS configurations.] Run multiple operating systems
      simultaneously, for instance for compatibility or QA purposes.
\item [Server consolidation.] Move multiple servers onto a single
      physical host with performance and fault isolation provided at
      virtual machine boundaries. 
\item [Cluster computing.] Management at VM granularity provides more
      flexibility than separately managing each physical host, but
      better control and isolation than single-system image solutions, 
      particularly by using live migration for load balancing. 
\item [Hardware support for custom OSes.] Allow development of new OSes
      while benefiting from the wide-ranging hardware support of
      existing OSes such as Linux.
\end{description}

\section{Structure of a Xen-Based System}

A Xen system has multiple layers, the lowest and most privileged of
which is Xen itself. 
Xen in turn may host multiple {\em guest} operating systems, each of
which is executed within a secure virtual machine (in Xen terminology,
a {\em domain}). Domains are scheduled by Xen to make effective use of
the available physical CPUs.  Each guest OS manages its own
applications, which includes responsibility for scheduling each
application within the time allotted to the VM by Xen.

The first domain, {\em domain 0}, is created automatically when the
system boots and has special management privileges. Domain 0 builds
other domains and manages their virtual devices. It also performs
administrative tasks such as suspending, resuming and migrating other
virtual machines.

Within domain 0, a process called \emph{xend} runs to manage the system.
\Xend is responsible for managing virtual machines and providing access
to their consoles.  Commands are issued to \xend over an HTTP
interface, either from a command-line tool or from a web browser.

\section{Hardware Support}

Xen currently runs only on the x86 architecture, requiring a `P6' or
newer processor (e.g. Pentium Pro, Celeron, Pentium II, Pentium III,
Pentium IV, Xeon, AMD Athlon, AMD Duron).  Multiprocessor machines are
supported, and we also have basic support for HyperThreading (SMT),
although this remains a topic for ongoing research. A port
specifically for x86/64 is in progress, although Xen already runs on
such systems in 32-bit legacy mode. In addition a port to the IA64
architecture is approaching completion. We hope to add other
architectures such as PPC and ARM in due course.


Xen can currently use up to 4GB of memory.  It is possible for x86
machines to address up to 64GB of physical memory but there are no
current plans to support these systems: The x86/64 port is the
planned route to supporting larger memory sizes.

Xen offloads most of the hardware support issues to the guest OS
running in Domain~0.  Xen itself contains only the code required to
detect and start secondary processors, set up interrupt routing, and
perform PCI bus enumeration.  Device drivers run within a privileged
guest OS rather than within Xen itself. This approach provides
compatibility with the majority of device hardware supported by Linux.
The default XenLinux build contains support for relatively modern
server-class network and disk hardware, but you can add support for
other hardware by configuring your XenLinux kernel in the normal way.

\section{History}

Xen was originally developed by the Systems Research Group at the
University of Cambridge Computer Laboratory as part of the XenoServers
project, funded by the UK-EPSRC.
XenoServers aim to provide a `public infrastructure for
global distributed computing', and Xen plays a key part in that,
allowing us to efficiently partition a single machine to enable
multiple independent clients to run their operating systems and
applications in an environment providing protection, resource
isolation and accounting.  The project web page contains further
information along with pointers to papers and technical reports:
\path{http://www.cl.cam.ac.uk/xeno} 

Xen has since grown into a fully-fledged project in its own right,
enabling us to investigate interesting research issues regarding the
best techniques for virtualising resources such as the CPU, memory,
disk and network.  The project has been bolstered by support from
Intel Research Cambridge, and HP Labs, who are now working closely
with us.

Xen was first described in a paper presented at SOSP in
2003\footnote{\tt
http://www.cl.cam.ac.uk/netos/papers/2003-xensosp.pdf}, and the first
public release (1.0) was made that October.  Since then, Xen has
significantly matured and is now used in production scenarios on
many sites.

Xen 2.0 features greatly enhanced hardware support, configuration
flexibility, usability and a larger complement of supported operating
systems. This latest release takes Xen a step closer to becoming the 
definitive open source solution for virtualisation.

\chapter{Installation}

The Xen distribution includes three main components: Xen itself, ports
of Linux 2.4 and 2.6 and NetBSD to run on Xen, and the user-space
tools required to manage a Xen-based system.  This chapter describes
how to install the Xen 2.0 distribution from source.  Alternatively,
there may be pre-built packages available as part of your operating
system distribution.

\section{Prerequisites}
\label{sec:prerequisites}

The following is a full list of prerequisites.  Items marked `$\dag$'
are required by the \xend control tools, and hence required if you
want to run more than one virtual machine; items marked `$*$' are only
required if you wish to build from source.
\begin{itemize}
\item A working Linux distribution using the GRUB bootloader and
running on a P6-class (or newer) CPU.
\item [$\dag$] The \path{iproute2} package. 
\item [$\dag$] The Linux bridge-utils\footnote{Available from 
{\tt http://bridge.sourceforge.net}} (e.g., \path{/sbin/brctl})
\item [$\dag$] An installation of Twisted v1.3 or
above\footnote{Available from {\tt
http://www.twistedmatrix.com}}. There may be a binary package
available for your distribution; alternatively it can be installed by
running `{\sl make install-twisted}' in the root of the Xen source
tree.
\item [$*$] Build tools (gcc v3.2.x or v3.3.x, binutils, GNU make).
\item [$*$] Development installation of libcurl (e.g., libcurl-devel) 
\item [$*$] Development installation of zlib (e.g., zlib-dev).
\item [$*$] Development installation of Python v2.2 or later (e.g., python-dev).
\item [$*$] \LaTeX, transfig and tgif are required to build the documentation.
\end{itemize}

Once you have satisfied the relevant prerequisites, you can 
now install either a binary or source distribution of Xen. 

\section{Installing from Binary Tarball} 

Pre-built tarballs are available for download from the Xen 
download page
\begin{quote} 
{\tt http://xen.sf.net}
\end{quote} 

Once you've downloaded the tarball, simply unpack and install: 
\begin{verbatim}
# tar zxvf xen-2.0-install.tgz
# cd xen-2.0-install
# sh ./install.sh 
\end{verbatim} 

Once you've installed the binaries you need to configure
your system as described in Section~\ref{s:configure}. 

\section{Installing from Source} 

This section describes how to obtain, build, and install 
Xen from source. 

\subsection{Obtaining the Source} 

The Xen source tree is available as either a compressed source tar
ball or as a clone of our master BitKeeper repository.

\begin{description} 
\item[Obtaining the Source Tarball]\mbox{} \\  
Stable versions (and daily snapshots) of the Xen source tree are
available as compressed tarballs from the Xen download page
\begin{quote} 
{\tt http://xen.sf.net}
\end{quote} 

\item[Using BitKeeper]\mbox{} \\  
If you wish to install Xen from a clone of our latest BitKeeper
repository then you will need to install the BitKeeper tools.
Download instructions for BitKeeper can be obtained by filling out the
form at:

\begin{quote} 
{\tt http://www.bitmover.com/cgi-bin/download.cgi}
\end{quote}
The public master BK repository for the 2.0 release lives at: 
\begin{quote}
{\tt bk://xen.bkbits.net/xen-2.0.bk}  
\end{quote} 
You can use BitKeeper to
download it and keep it updated with the latest features and fixes.

Change to the directory in which you want to put the source code, then
run:
\begin{verbatim}
# bk clone bk://xen.bkbits.net/xen-2.0.bk
\end{verbatim}

Under your current directory, a new directory named \path{xen-2.0.bk}
has been created, which contains all the source code for Xen, the OS
ports, and the control tools. You can update your repository with the
latest changes at any time by running:
\begin{verbatim}
# cd xen-2.0.bk # to change into the local repository
# bk pull       # to update the repository
\end{verbatim}
\end{description} 

%\section{The distribution}
%
%The Xen source code repository is structured as follows:
%
%\begin{description}
%\item[\path{tools/}] Xen node controller daemon (Xend), command line tools, 
%  control libraries
%\item[\path{xen/}] The Xen VMM.
%\item[\path{linux-*-xen-sparse/}] Xen support for Linux.
%\item[\path{linux-*-patches/}] Experimental patches for Linux.
%\item[\path{netbsd-*-xen-sparse/}] Xen support for NetBSD.
%\item[\path{docs/}] Various documentation files for users and developers.
%\item[\path{extras/}] Bonus extras.
%\end{description}

\subsection{Building from Source} 

The top-level Xen Makefile includes a target `world' that will do the
following:

\begin{itemize}
\item Build Xen
\item Build the control tools, including \xend
\item Download (if necessary) and unpack the Linux 2.6 source code,
      and patch it for use with Xen
\item Build a Linux kernel to use in domain 0 and a smaller
      unprivileged kernel, which can optionally be used for
      unprivileged virtual machines.
\end{itemize}


After the build has completed you should have a top-level 
directory called \path{dist/} in which all resulting targets 
will be placed; of particular interest are the two kernels 
XenLinux kernel images, one with a `-xen0' extension
which contains hardware device drivers and drivers for Xen's virtual
devices, and one with a `-xenU' extension that just contains the
virtual ones. These are found in \path{dist/install/boot/} along
with the image for Xen itself and the configuration files used
during the build. 

The NetBSD port can be built using: 
\begin{quote}
\begin{verbatim}
# make netbsd20
\end{verbatim} 
\end{quote} 
NetBSD port is built using a snapshot of the netbsd-2-0 cvs branch.
The snapshot is downloaded as part of the build process, if it is not
yet present in the \path{NETBSD\_SRC\_PATH} search path.  The build
process also downloads a toolchain which includes all the tools
necessary to build the NetBSD kernel under Linux.

To customize further the set of kernels built you need to edit
the top-level Makefile. Look for the line: 

\begin{quote}
\begin{verbatim}
KERNELS ?= mk.linux-2.6-xen0 mk.linux-2.6-xenU
\end{verbatim} 
\end{quote} 

You can edit this line to include any set of operating system kernels
which have configurations in the top-level \path{buildconfigs/}
directory, for example \path{mk.linux-2.4-xenU} to build a Linux 2.4
kernel containing only virtual device drivers.

%% Inspect the Makefile if you want to see what goes on during a build.
%% Building Xen and the tools is straightforward, but XenLinux is more
%% complicated.  The makefile needs a `pristine' Linux kernel tree to which
%% it will then add the Xen architecture files.  You can tell the
%% makefile the location of the appropriate Linux compressed tar file by
%% setting the LINUX\_SRC environment variable, e.g. \\
%% \verb!# LINUX_SRC=/tmp/linux-2.6.9.tar.bz2 make world! \\ or by
%% placing the tar file somewhere in the search path of {\tt
%% LINUX\_SRC\_PATH} which defaults to `{\tt .:..}'.  If the makefile
%% can't find a suitable kernel tar file it attempts to download it from
%% kernel.org (this won't work if you're behind a firewall).

%% After untaring the pristine kernel tree, the makefile uses the {\tt
%% mkbuildtree} script to add the Xen patches to the kernel. 


%% The procedure is similar to build the Linux 2.4 port: \\
%% \verb!# LINUX_SRC=/path/to/linux2.4/source make linux24!


%% \framebox{\parbox{5in}{
%% {\bf Distro specific:} \\
%% {\it Gentoo} --- if not using udev (most installations, currently), you'll need
%% to enable devfs and devfs mount at boot time in the xen0 config.
%% }}

\subsection{Custom XenLinux Builds}

% If you have an SMP machine you may wish to give the {\tt '-j4'}
% argument to make to get a parallel build.

If you wish to build a customized XenLinux kernel (e.g. to support
additional devices or enable distribution-required features), you can
use the standard Linux configuration mechanisms, specifying that the
architecture being built for is \path{xen}, e.g:
\begin{quote}
\begin{verbatim} 
# cd linux-2.6.9-xen0 
# make ARCH=xen xconfig 
# cd ..
# make
\end{verbatim} 
\end{quote} 

You can also copy an existing Linux configuration (\path{.config}) 
into \path{linux-2.6.9-xen0} and execute:  
\begin{quote}
\begin{verbatim} 
# make ARCH=xen oldconfig 
\end{verbatim} 
\end{quote} 

You may be prompted with some Xen-specific options; we 
advise accepting the defaults for these options.

Note that the only difference between the two types of Linux kernel
that are built is the configuration file used for each.  The "U"
suffixed (unprivileged) versions don't contain any of the physical
hardware device drivers, leading to a 30\% reduction in size; hence
you may prefer these for your non-privileged domains.  The `0'
suffixed privileged versions can be used to boot the system, as well
as in driver domains and unprivileged domains.


\subsection{Installing the Binaries}


The files produced by the build process are stored under the
\path{dist/install/} directory. To install them in their default
locations, do:
\begin{quote}
\begin{verbatim}
# make install
\end{verbatim} 
\end{quote}


Alternatively, users with special installation requirements may wish
to install them manually by copying the files to their appropriate
destinations.

%% Files in \path{install/boot/} include:
%% \begin{itemize}
%% \item \path{install/boot/xen.gz} The Xen 'kernel'
%% \item \path{install/boot/vmlinuz-2.6.9-xen0}  Domain 0 XenLinux kernel
%% \item \path{install/boot/vmlinuz-2.6.9-xenU}  Unprivileged XenLinux kernel
%% \end{itemize}

The \path{dist/install/boot} directory will also contain the config files
used for building the XenLinux kernels, and also versions of Xen and
XenLinux kernels that contain debug symbols (\path{xen-syms} and
\path{vmlinux-syms-2.6.9-xen0}) which are essential for interpreting crash
dumps.  Retain these files as the developers may wish to see them if
you post on the mailing list.





\section{Configuration}
\label{s:configure}
Once you have built and installed the Xen distribution, it is 
simple to prepare the machine for booting and running Xen. 

\subsection{GRUB Configuration}

An entry should be added to \path{grub.conf} (often found under
\path{/boot/} or \path{/boot/grub/}) to allow Xen / XenLinux to boot.
This file is sometimes called \path{menu.lst}, depending on your
distribution.  The entry should look something like the following:

{\small
\begin{verbatim}
title Xen 2.0 / XenLinux 2.6.9
  kernel /boot/xen.gz dom0_mem=131072
  module /boot/vmlinuz-2.6.9-xen0 root=/dev/sda4 ro console=tty0
\end{verbatim}
}

The kernel line tells GRUB where to find Xen itself and what boot
parameters should be passed to it (in this case, setting domain 0's
memory allocation and the settings for the serial port). For more
details on the various Xen boot parameters see Section~\ref{s:xboot}. 

The module line of the configuration describes the location of the
XenLinux kernel that Xen should start and the parameters that should
be passed to it (these are standard Linux parameters, identifying the
root device and specifying it be initially mounted read only and
instructing that console output be sent to the screen).  Some
distributions such as SuSE do not require the \path{ro} parameter.

%% \framebox{\parbox{5in}{
%% {\bf Distro specific:} \\
%% {\it SuSE} --- Omit the {\tt ro} option from the XenLinux kernel
%% command line, since the partition won't be remounted rw during boot.
%% }}


If you want to use an initrd, just add another \path{module} line to
the configuration, as usual:
{\small
\begin{verbatim}
  module /boot/my_initrd.gz
\end{verbatim}
}

As always when installing a new kernel, it is recommended that you do
not delete existing menu options from \path{menu.lst} --- you may want
to boot your old Linux kernel in future, particularly if you
have problems.


\subsection{Serial Console (optional)}

%%   kernel /boot/xen.gz dom0_mem=131072 com1=115200,8n1
%%   module /boot/vmlinuz-2.6.9-xen0 root=/dev/sda4 ro 


In order to configure Xen serial console output, it is necessary to add 
an boot option to your GRUB config; e.g. replace the above kernel line 
with: 
\begin{quote}
{\small
\begin{verbatim}
   kernel /boot/xen.gz dom0_mem=131072 com1=115200,8n1
\end{verbatim}}
\end{quote} 

This configures Xen to output on COM1 at 115,200 baud, 8 data bits, 
1 stop bit and no parity. Modify these parameters for your set up. 

One can also configure XenLinux to share the serial console; to 
achieve this append ``\path{console=ttyS0}'' to your 
module line. 


If you wish to be able to log in over the XenLinux serial console it
is necessary to add a line into \path{/etc/inittab}, just as per 
regular Linux. Simply add the line:
\begin{quote}
{\small 
{\tt c:2345:respawn:/sbin/mingetty ttyS0}
}
\end{quote} 

and you should be able to log in. Note that to successfully log in 
as root over the serial line will require adding \path{ttyS0} to
\path{/etc/securetty} in most modern distributions. 

\subsection{TLS Libraries}

Users of the XenLinux 2.6 kernel should disable Thread Local Storage
(e.g.\ by doing a \path{mv /lib/tls /lib/tls.disabled}) before
attempting to run with a XenLinux kernel\footnote{If you boot without first
disabling TLS, you will get a warning message during the boot
process. In this case, simply perform the rename after the machine is
up and then run \texttt{/sbin/ldconfig} to make it take effect.}.  You can
always reenable it by restoring the directory to its original location
(i.e.\ \path{mv /lib/tls.disabled /lib/tls}).

The reason for this is that the current TLS implementation uses
segmentation in a way that is not permissible under Xen.  If TLS is
not disabled, an emulation mode is used within Xen which reduces
performance substantially.

We hope that this issue can be resolved by working with Linux
distribution vendors to implement a minor backward-compatible change
to the TLS library.

\section{Booting Xen} 

It should now be possible to restart the system and use Xen.  Reboot
as usual but choose the new Xen option when the Grub screen appears.

What follows should look much like a conventional Linux boot.  The
first portion of the output comes from Xen itself, supplying low level
information about itself and the machine it is running on.  The
following portion of the output comes from XenLinux.

You may see some errors during the XenLinux boot.  These are not
necessarily anything to worry about --- they may result from kernel
configuration differences between your XenLinux kernel and the one you
usually use.

When the boot completes, you should be able to log into your system as
usual.  If you are unable to log in to your system running Xen, you
should still be able to reboot with your normal Linux kernel.


\chapter{Starting Additional Domains}

The first step in creating a new domain is to prepare a root
filesystem for it to boot off.  Typically, this might be stored in a
normal partition, an LVM or other volume manager partition, a disk
file or on an NFS server.  A simple way to do this is simply to boot
from your standard OS install CD and install the distribution into
another partition on your hard drive.

To start the \xend control daemon, type
\begin{quote}
\verb!# xend start!
\end{quote}
If you
wish the daemon to start automatically, see the instructions in
Section~\ref{s:xend}. Once the daemon is running, you can use the
\path{xm} tool to monitor and maintain the domains running on your
system. This chapter provides only a brief tutorial: we provide full
details of the \path{xm} tool in the next chapter. 

%\section{From the web interface}
%
%Boot the Xen machine and start Xensv (see Chapter~\ref{cha:xensv} for
%more details) using the command: \\
%\verb_# xensv start_ \\
%This will also start Xend (see Chapter~\ref{cha:xend} for more information).
%
%The domain management interface will then be available at {\tt
%http://your\_machine:8080/}.  This provides a user friendly wizard for
%starting domains and functions for managing running domains.
%
%\section{From the command line}


\section{Creating a Domain Configuration File} 

Before you can start an additional domain, you must create a
configuration file. We provide two example files which you 
can use as a starting point: 
\begin{itemize} 
  \item \path{/etc/xen/xmexample1} is a simple template configuration file
    for describing a single VM.

  \item \path{/etc/xen/xmexample2} file is a template description that
    is intended to be reused for multiple virtual machines.  Setting
    the value of the \path{vmid} variable on the \path{xm} command line
    fills in parts of this template.
\end{itemize} 

Copy one of these files and edit it as appropriate.
Typical values you may wish to edit include: 

\begin{quote}
\begin{description}
\item[kernel] Set this to the path of the kernel you compiled for use
              with Xen (e.g.\  \path{kernel = '/boot/vmlinuz-2.6.9-xenU'})
\item[memory] Set this to the size of the domain's memory in
megabytes (e.g.\ \path{memory = 64})
\item[disk] Set the first entry in this list to calculate the offset
of the domain's root partition, based on the domain ID.  Set the
second to the location of \path{/usr} if you are sharing it between
domains (e.g.\ \path{disk = ['phy:your\_hard\_drive\%d,sda1,w' \%
(base\_partition\_number + vmid), 'phy:your\_usr\_partition,sda6,r' ]}
\item[dhcp] Uncomment the dhcp variable, so that the domain will
receive its IP address from a DHCP server (e.g.\ \path{dhcp='dhcp'})
\end{description}
\end{quote}

You may also want to edit the {\bf vif} variable in order to choose
the MAC address of the virtual ethernet interface yourself.  For
example: 
\begin{quote}
\verb_vif = ['mac=00:06:AA:F6:BB:B3']_
\end{quote}
If you do not set this variable, \xend will automatically generate a
random MAC address from an unused range.


\section{Booting the Domain}

The \path{xm} tool provides a variety of commands for managing domains.
Use the \path{create} command to start new domains. Assuming you've 
created a configuration file \path{myvmconf} based around
\path{/etc/xen/xmexample2}, to start a domain with virtual 
machine ID~1 you should type: 

\begin{quote}
\begin{verbatim}
# xm create -c myvmconf vmid=1
\end{verbatim}
\end{quote}


The \path{-c} switch causes \path{xm} to turn into the domain's
console after creation.  The \path{vmid=1} sets the \path{vmid}
variable used in the \path{myvmconf} file. 


You should see the console boot messages from the new domain 
appearing in the terminal in which you typed the command, 
culminating in a login prompt. 


\section{Example: ttylinux}

Ttylinux is a very small Linux distribution, designed to require very
few resources.  We will use it as a concrete example of how to start a
Xen domain.  Most users will probably want to install a full-featured
distribution once they have mastered the basics.

\begin{enumerate}
\item Download and extract the ttylinux disk image from the Files
section of the project's SourceForge site (see 
\path{http://sf.net/projects/xen/}).
\item Create a configuration file like the following:
\begin{verbatim}
kernel = "/boot/vmlinuz-2.6.9-xenU"
memory = 64
name = "ttylinux"
nics = 1
ip = "1.2.3.4"
disk = ['file:/path/to/ttylinux/rootfs,sda1,w']
root = "/dev/sda1 ro"
\end{verbatim}
\item Now start the domain and connect to its console:
\begin{verbatim}
xm create configfile -c
\end{verbatim}
\item Login as root, password root.
\end{enumerate}


\section{Starting / Stopping Domains Automatically}

It is possible to have certain domains start automatically at boot
time and to have dom0 wait for all running domains to shutdown before
it shuts down the system.

To specify a domain is to start at boot-time, place its
configuration file (or a link to it) under \path{/etc/xen/auto/}.

A Sys-V style init script for RedHat and LSB-compliant systems is
provided and will be automatically copied to \path{/etc/init.d/}
during install.  You can then enable it in the appropriate way for
your distribution.

For instance, on RedHat:

\begin{quote}
\verb_# chkconfig --add xendomains_
\end{quote}

By default, this will start the boot-time domains in runlevels 3, 4
and 5.

You can also use the \path{service} command to run this script
manually, e.g:

\begin{quote}
\verb_# service xendomains start_

Starts all the domains with config files under /etc/xen/auto/.
\end{quote}


\begin{quote}
\verb_# service xendomains stop_

Shuts down ALL running Xen domains.
\end{quote}

\chapter{Domain Management Tools}

The previous chapter described a simple example of how to configure
and start a domain.  This chapter summarises the tools available to
manage running domains.

\section{Command-line Management}

Command line management tasks are also performed using the \path{xm}
tool.  For online help for the commands available, type:
\begin{quote}
\verb_# xm help_
\end{quote}

You can also type \path{xm help $<$command$>$} for more information 
on a given command. 

\subsection{Basic Management Commands}

The most important \path{xm} commands are: 
\begin{quote}
\verb_# xm list_: Lists all domains running.\\
\verb_# xm consoles_ : Gives information about the domain consoles.\\
\verb_# xm console_: Opens a console to a domain (e.g.\
  \verb_# xm console myVM_
\end{quote}

\subsection{\tt xm list}

The output of \path{xm list} is in rows of the following format:
\begin{center}
{\tt name domid memory cpu state cputime console}
\end{center}

\begin{quote}
\begin{description}
\item[name]  The descriptive name of the virtual machine.
\item[domid] The number of the domain ID this virtual machine is running in.
\item[memory] Memory size in megabytes.
\item[cpu]   The CPU this domain is running on.
\item[state] Domain state consists of 5 fields:
  \begin{description}
  \item[r] running
  \item[b] blocked
  \item[p] paused
  \item[s] shutdown
  \item[c] crashed
  \end{description}
\item[cputime] How much CPU time (in seconds) the domain has used so far.
\item[console] TCP port accepting connections to the domain's console.
\end{description}
\end{quote}

The \path{xm list} command also supports a long output format when the
\path{-l} switch is used.  This outputs the fulls details of the
running domains in \xend's SXP configuration format.

For example, suppose the system is running the ttylinux domain as
described earlier.  The list command should produce output somewhat
like the following:
\begin{verbatim}
# xm list
Name              Id  Mem(MB)  CPU  State  Time(s)  Console
Domain-0           0      251    0  r----    172.2        
ttylinux           5       63    0  -b---      3.0    9605
\end{verbatim}

Here we can see the details for the ttylinux domain, as well as for
domain 0 (which, of course, is always running).  Note that the console
port for the ttylinux domain is 9605.  This can be connected to by TCP
using a terminal program (e.g. \path{telnet} or, better, 
\path{xencons}).  The simplest way to connect is to use the \path{xm console}
command, specifying the domain name or ID.  To connect to the console
of the ttylinux domain, we could use any of the following: 
\begin{verbatim}
# xm console ttylinux
# xm console 5
# xencons localhost 9605
\end{verbatim}

\section{Domain Save and Restore}

The administrator of a Xen system may suspend a virtual machine's
current state into a disk file in domain 0, allowing it to be resumed
at a later time.

The ttylinux domain described earlier can be suspended to disk using
the command:
\begin{verbatim}
# xm save ttylinux ttylinux.xen
\end{verbatim}

This will stop the domain named `ttylinux' and save its current state
into a file called \path{ttylinux.xen}.

To resume execution of this domain, use the \path{xm restore} command:
\begin{verbatim}
# xm restore ttylinux.xen
\end{verbatim}

This will restore the state of the domain and restart it.  The domain
will carry on as before and the console may be reconnected using the
\path{xm console} command, as above.

\section{Live Migration}

Live migration is used to transfer a domain between physical hosts
whilst that domain continues to perform its usual activities --- from
the user's perspective, the migration should be imperceptible.

To perform a live migration, both hosts must be running Xen / \xend and
the destination host must have sufficient resources (e.g. memory
capacity) to accommodate the domain after the move. Furthermore we
currently require both source and destination machines to be on the 
same L2 subnet. 

Currently, there is no support for providing automatic remote access
to filesystems stored on local disk when a domain is migrated.
Administrators should choose an appropriate storage solution
(i.e. SAN, NAS, etc.) to ensure that domain filesystems are also
available on their destination node. GNBD is a good method for
exporting a volume from one machine to another. iSCSI can do a similar
job, but is more complex to set up.

When a domain migrates, it's MAC and IP address move with it, thus it
is only possible to migrate VMs within the same layer-2 network and IP
subnet. If the destination node is on a different subnet, the
administrator would need to manually configure a suitable etherip or
IP tunnel in the domain 0 of the remote node. 

A domain may be migrated using the \path{xm migrate} command.  To
live migrate a domain to another machine, we would use
the command:

\begin{verbatim}
# xm migrate --live mydomain destination.ournetwork.com
\end{verbatim}

Without the \path{--live} flag, \xend simply stops the domain and
copies the memory image over to the new node and restarts it. Since
domains can have large allocations this can be quite time consuming,
even on a Gigabit network. With the \path{--live} flag \xend attempts
to keep the domain running while the migration is in progress,
resulting in typical `downtimes' of just 60--300ms.

For now it will be necessary to reconnect to the domain's console on
the new machine using the \path{xm console} command.  If a migrated
domain has any open network connections then they will be preserved,
so SSH connections do not have this limitation.

\section{Managing Domain Memory}

XenLinux domains have the ability to relinquish / reclaim machine
memory at the request of the administrator or the user of the domain.

\subsection{Setting memory footprints from dom0}

The machine administrator can request that a domain alter its memory
footprint using the \path{xm balloon} command.  For instance, we can
request that our example ttylinux domain reduce its memory footprint
to 32 megabytes.

\begin{verbatim}
# xm balloon ttylinux 32
\end{verbatim}

We can now see the result of this in the output of \path{xm list}:

\begin{verbatim}
# xm list
Name              Id  Mem(MB)  CPU  State  Time(s)  Console
Domain-0           0      251    0  r----    172.2        
ttylinux           5       31    0  -b---      4.3    9605
\end{verbatim}

The domain has responded to the request by returning memory to Xen. We
can restore the domain to its original size using the command line:

\begin{verbatim}
# xm balloon ttylinux 64
\end{verbatim}

\subsection{Setting memory footprints from within a domain}

The virtual file \path{/proc/xen/memory\_target} allows the owner of a
domain to adjust their own memory footprint.  Reading the file
(e.g. \path{cat /proc/xen/memory\_target}) prints out the current
memory footprint of the domain.  Writing the file
(e.g. \path{echo new\_target > /proc/xen/memory\_target}) requests
that the kernel adjust the domain's memory footprint to a new value.

\subsection{Setting memory limits}

Xen associates a memory size limit with each domain.  By default, this
is the amount of memory the domain is originally started with,
preventing the domain from ever growing beyond this size.  To permit a
domain to grow beyond its original allocation or to prevent a domain
you've shrunk from reclaiming the memory it relinquished, use the 
\path{xm maxmem} command.

\chapter{Domain Filesystem Storage}

It is possible to directly export any Linux block device in dom0 to
another domain, or to export filesystems / devices to virtual machines
using standard network protocols (e.g. NBD, iSCSI, NFS, etc).  This
chapter covers some of the possibilities.


\section{Exporting Physical Devices as VBDs} 

One of the simplest configurations is to directly export 
individual partitions from domain 0 to other domains. To 
achieve this use the \path{phy:} specifier in your domain 
configuration file. For example a line like
\begin{quote}
\verb_disk = ['phy:hda3,sda1,w']_
\end{quote}
specifies that the partition \path{/dev/hda3} in domain 0 
should be exported read-write to the new domain as \path{/dev/sda1}; 
one could equally well export it as \path{/dev/hda} or 
\path{/dev/sdb5} should one wish. 

In addition to local disks and partitions, it is possible to export
any device that Linux considers to be ``a disk'' in the same manner.
For example, if you have iSCSI disks or GNBD volumes imported into
domain 0 you can export these to other domains using the \path{phy:}
disk syntax. E.g.:
\begin{quote}
\verb_disk = ['phy:vg/lvm1,sda2,w']_
\end{quote}



\begin{center}
\framebox{\bf Warning: Block device sharing}
\end{center}
\begin{quote}
Block devices should typically only be shared between domains in a
read-only fashion otherwise the Linux kernel's file systems will get
very confused as the file system structure may change underneath them
(having the same ext3 partition mounted rw twice is a sure fire way to
cause irreparable damage)!  \Xend will attempt to prevent you from
doing this by checking that the device is not mounted read-write in
domain 0, and hasn't already been exported read-write to another
domain.
If you want read-write sharing, export the directory to other domains
via NFS from domain0 (or use a cluster file system such as GFS or
ocfs2).

\end{quote}


\section{Using File-backed VBDs}

It is also possible to use a file in Domain 0 as the primary storage
for a virtual machine.  As well as being convenient, this also has the
advantage that the virtual block device will be {\em sparse} --- space
will only really be allocated as parts of the file are used.  So if a
virtual machine uses only half of its disk space then the file really
takes up half of the size allocated.

For example, to create a 2GB sparse file-backed virtual block device
(actually only consumes 1KB of disk):
\begin{quote}
\verb_# dd if=/dev/zero of=vm1disk bs=1k seek=2048k count=1_
\end{quote}

Make a file system in the disk file: 
\begin{quote}
\verb_# mkfs -t ext3 vm1disk_
\end{quote}

(when the tool asks for confirmation, answer `y')

Populate the file system e.g. by copying from the current root:
\begin{quote}
\begin{verbatim}
# mount -o loop vm1disk /mnt
# cp -ax /{root,dev,var,etc,usr,bin,sbin,lib} /mnt
# mkdir /mnt/{proc,sys,home,tmp}
\end{verbatim}
\end{quote}

Tailor the file system by editing \path{/etc/fstab},
\path{/etc/hostname}, etc (don't forget to edit the files in the
mounted file system, instead of your domain 0 filesystem, e.g. you
would edit \path{/mnt/etc/fstab} instead of \path{/etc/fstab} ).  For
this example put \path{/dev/sda1} to root in fstab.

Now unmount (this is important!):
\begin{quote}
\verb_# umount /mnt_
\end{quote}

In the configuration file set:
\begin{quote}
\verb_disk = ['file:/full/path/to/vm1disk,sda1,w']_
\end{quote}

As the virtual machine writes to its `disk', the sparse file will be
filled in and consume more space up to the original 2GB.


\section{Using LVM-backed VBDs}

A particularly appealing solution is to use LVM volumes 
as backing for domain file-systems since this allows dynamic
growing/shrinking of volumes as well as snapshot and other 
features. 

To initialise a partition to support LVM volumes:
\begin{quote}
\begin{verbatim} 
# pvcreate /dev/sda10		
\end{verbatim} 
\end{quote}

Create a volume group named `vg' on the physical partition:
\begin{quote}
\begin{verbatim} 
# vgcreate vg /dev/sda10
\end{verbatim} 
\end{quote}

Create a logical volume of size 4GB named `myvmdisk1':
\begin{quote}
\begin{verbatim} 
# lvcreate -L4096M -n myvmdisk1 vg
\end{verbatim} 
\end{quote}

You should now see that you have a \path{/dev/vg/myvmdisk1}
Make a filesystem, mount it and populate it, e.g.:
\begin{quote}
\begin{verbatim} 
# mkfs -t ext3 /dev/vg/myvmdisk1
# mount /dev/vg/myvmdisk1 /mnt
# cp -ax / /mnt
# umount /mnt
\end{verbatim} 
\end{quote}

Now configure your VM with the following disk configuration:
\begin{quote}
\begin{verbatim} 
 disk = [ 'phy:vg/myvmdisk1,sda1,w' ]
\end{verbatim} 
\end{quote}

LVM enables you to grow the size of logical volumes, but you'll need
to resize the corresponding file system to make use of the new
space. Some file systems (e.g. ext3) now support on-line resize.  See
the LVM manuals for more details.

You can also use LVM for creating copy-on-write clones of LVM
volumes (known as writable persistent snapshots in LVM
terminology). This facility is new in Linux 2.6.8, so isn't as
stable as one might hope. In particular, using lots of CoW LVM
disks consumes a lot of dom0 memory, and error conditions such as
running out of disk space are not handled well. Hopefully this
will improve in future.

To create two copy-on-write clone of the above file system you
would use the following commands:

\begin{quote}
\begin{verbatim} 
# lvcreate -s -L1024M -n myclonedisk1 /dev/vg/myvmdisk1
# lvcreate -s -L1024M -n myclonedisk2 /dev/vg/myvmdisk1
\end{verbatim} 
\end{quote}

Each of these can grow to have 1GB of differences from the master
volume. You can grow the amount of space for storing the
differences using the lvextend command, e.g.:
\begin{quote}
\begin{verbatim} 
# lvextend +100M /dev/vg/myclonedisk1
\end{verbatim} 
\end{quote}

Don't let the `differences volume' ever fill up otherwise LVM gets
rather confused. It may be possible to automate the growing
process by using \path{dmsetup wait} to spot the volume getting full
and then issue an \path{lvextend}.

In principle, it is possible to continue writing to the volume
that has been cloned (the changes will not be visible to the
clones), but we wouldn't recommend this: have the cloned volume
as a `pristine' file system install that isn't mounted directly
by any of the virtual machines.


\section{Using NFS Root}

First, populate a root filesystem in a directory on the server
machine. This can be on a distinct physical machine, or simply 
run within a virtual machine on the same node.

Now configure the NFS server to export this filesystem over the
network by adding a line to \path{/etc/exports}, for instance:

\begin{quote}
\begin{small}
\begin{verbatim}
/export/vm1root      1.2.3.4/24 (rw,sync,no_root_squash)
\end{verbatim}
\end{small}
\end{quote}

Finally, configure the domain to use NFS root.  In addition to the
normal variables, you should make sure to set the following values in
the domain's configuration file:

\begin{quote}
\begin{small}
\begin{verbatim}
root       = '/dev/nfs'
nfs_server = '2.3.4.5'       # substitute IP address of server 
nfs_root   = '/path/to/root' # path to root FS on the server
\end{verbatim}
\end{small}
\end{quote}

The domain will need network access at boot time, so either statically
configure an IP address (Using the config variables \path{ip}, 
\path{netmask}, \path{gateway}, \path{hostname}) or enable DHCP (
\path{dhcp='dhcp'}).

Note that the Linux NFS root implementation is known to have stability
problems under high load (this is not a Xen-specific problem), so this
configuration may not be appropriate for critical servers.


\part{User Reference Documentation}

\chapter{Control Software} 

The Xen control software includes the \xend node control daemon (which 
must be running), the xm command line tools, and the prototype 
xensv web interface. 

\section{\Xend (node control daemon)}
\label{s:xend}

The Xen Daemon (\Xend) performs system management functions related to
virtual machines.  It forms a central point of control for a machine
and can be controlled using an HTTP-based protocol.  \Xend must be
running in order to start and manage virtual machines.

\Xend must be run as root because it needs access to privileged system
management functions.  A small set of commands may be issued on the
\xend command line:

\begin{tabular}{ll}
\verb!# xend start! & start \xend, if not already running \\
\verb!# xend stop!  & stop \xend if already running       \\
\verb!# xend restart! & restart \xend if running, otherwise start it \\
% \verb!# xend trace_start! & start \xend, with very detailed debug logging \\
\verb!# xend status! & indicates \xend status by its return code
\end{tabular}

A SysV init script called {\tt xend} is provided to start \xend at boot
time.  {\tt make install} installs this script in {\path{/etc/init.d}.
To enable it, you have to make symbolic links in the appropriate
runlevel directories or use the {\tt chkconfig} tool, where available.

Once \xend is running, more sophisticated administration can be done
using the xm tool (see Section~\ref{s:xm}) and the experimental
Xensv web interface (see Section~\ref{s:xensv}).

As \xend runs, events will be logged to \path{/var/log/xend.log} and, 
if the migration assistant daemon (\path{xfrd}) has been started, 
\path{/var/log/xfrd.log}. These may be of use for troubleshooting
problems.

\section{Xm (command line interface)}
\label{s:xm}

The xm tool is the primary tool for managing Xen from the console.
The general format of an xm command line is:

\begin{verbatim}
# xm command [switches] [arguments] [variables]
\end{verbatim}

The available {\em switches} and {\em arguments} are dependent on the
{\em command} chosen.  The {\em variables} may be set using
declarations of the form {\tt variable=value} and command line
declarations override any of the values in the configuration file
being used, including the standard variables described above and any
custom variables (for instance, the \path{xmdefconfig} file uses a
{\tt vmid} variable).

The available commands are as follows:

\begin{description}
\item[balloon] Request a domain to adjust its memory footprint.
\item[create] Create a new domain.
\item[destroy] Kill a domain immediately.
\item[list] List running domains.
\item[shutdown] Ask a domain to shutdown.
\item[dmesg] Fetch the Xen (not Linux!) boot output.
\item[consoles] Lists the available consoles.
\item[console] Connect to the console for a domain.
\item[help] Get help on xm commands.
\item[save] Suspend a domain to disk.
\item[restore] Restore a domain from disk.
\item[pause] Pause a domain's execution.
\item[unpause] Unpause a domain.
\item[pincpu] Pin a domain to a CPU.
\item[bvt] Set BVT scheduler parameters for a domain.
\item[bvt\_ctxallow] Set the BVT context switching allowance for the system.
\item[atropos] Set the atropos parameters for a domain.
\item[rrobin] Set the round robin time slice for the system.
\item[info] Get information about the Xen host.
\item[call] Call a \xend HTTP API function directly.
\end{description}

For a detailed overview of switches, arguments and variables to each command
try
\begin{quote}
\begin{verbatim}
# xm help command
\end{verbatim}
\end{quote}

\section{Xensv (web control interface)}
\label{s:xensv}

Xensv is the experimental web control interface for managing a Xen
machine.  It can be used to perform some (but not yet all) of the
management tasks that can be done using the xm tool.

It can be started using:
\begin{quote}
\verb_# xensv start_
\end{quote}
and stopped using: 
\begin{quote}
\verb_# xensv stop_
\end{quote}

By default, Xensv will serve out the web interface on port 8080.  This
can be changed by editing 
\path{/usr/lib/python2.3/site-packages/xen/sv/params.py}.

Once Xensv is running, the web interface can be used to create and
manage running domains.




\chapter{Domain Configuration}
\label{cha:config}

The following contains the syntax of the domain configuration 
files and description of how to further specify networking, 
driver domain and general scheduling behaviour. 

\section{Configuration Files}
\label{s:cfiles}

Xen configuration files contain the following standard variables.
Unless otherwise stated, configuration items should be enclosed in
quotes: see \path{/etc/xen/xmexample1} and \path{/etc/xen/xmexample2} 
for concrete examples of the syntax.

\begin{description}
\item[kernel] Path to the kernel image 
\item[ramdisk] Path to a ramdisk image (optional).
% \item[builder] The name of the domain build function (e.g. {\tt'linux'} or {\tt'netbsd'}.
\item[memory] Memory size in megabytes.
\item[cpu] CPU to run this domain on, or {\tt -1} for
  auto-allocation. 
\item[console] Port to export the domain console on (default 9600 + domain ID).
\item[nics] Number of virtual network interfaces.
\item[vif] List of MAC addresses (random addresses are assigned if not
  given) and bridges to use for the domain's network interfaces, e.g.
\begin{verbatim}
vif = [ 'mac=aa:00:00:00:00:11, bridge=xen-br0',
        'bridge=xen-br1' ]
\end{verbatim}
  to assign a MAC address and bridge to the first interface and assign
  a different bridge to the second interface, leaving \xend to choose
  the MAC address.
\item[disk] List of block devices to export to the domain,  e.g. \\
  \verb_disk = [ 'phy:hda1,sda1,r' ]_ \\
  exports physical device \path{/dev/hda1} to the domain 
  as \path{/dev/sda1} with read-only access. Exporting a disk read-write 
  which is currently mounted is dangerous -- if you are \emph{certain}
  you wish to do this, you can specify \path{w!} as the mode. 
\item[dhcp] Set to {\tt 'dhcp'} if you want to use DHCP to configure
  networking. 
\item[netmask] Manually configured IP netmask.
\item[gateway] Manually configured IP gateway. 
\item[hostname] Set the hostname for the virtual machine.
\item[root] Specify the root device parameter on the kernel command
  line. 
\item[nfs\_server] IP address for the NFS server (if any). 
\item[nfs\_root] Path of the root filesystem on the NFS server (if any).
\item[extra] Extra string to append to the kernel command line (if
  any) 
\item[restart] Three possible options:
  \begin{description}
  \item[always] Always restart the domain, no matter what
                its exit code is.
  \item[never]  Never restart the domain.
  \item[onreboot] Restart the domain iff it requests reboot.
  \end{description}
\end{description}

For additional flexibility, it is also possible to include Python
scripting commands in configuration files.  An example of this is the
\path{xmexample2} file, which uses Python code to handle the 
\path{vmid} variable.


%\part{Advanced Topics}

\section{Network Configuration}

For many users, the default installation should work `out of the box'.
More complicated network setups, for instance with multiple ethernet
interfaces and/or existing bridging setups will require some
special configuration.

The purpose of this section is to describe the mechanisms provided by
\xend to allow a flexible configuration for Xen's virtual networking.

\subsection{Xen virtual network topology}

Each domain network interface is connected to a virtual network
interface in dom0 by a point to point link (effectively a `virtual
crossover cable').  These devices are named {\tt
vif$<$domid$>$.$<$vifid$>$} (e.g. {\tt vif1.0} for the first interface
in domain 1, {\tt vif3.1} for the second interface in domain 3).

Traffic on these virtual interfaces is handled in domain 0 using
standard Linux mechanisms for bridging, routing, rate limiting, etc.
Xend calls on two shell scripts to perform initial configuration of
the network and configuration of new virtual interfaces.  By default,
these scripts configure a single bridge for all the virtual
interfaces.  Arbitrary routing / bridging configurations can be
configured by customising the scripts, as described in the following
section.

\subsection{Xen networking scripts}

Xen's virtual networking is configured by two shell scripts (by
default \path{network} and \path{vif-bridge}).  These are
called automatically by \xend when certain events occur, with
arguments to the scripts providing further contextual information.
These scripts are found by default in \path{/etc/xen/scripts}.  The
names and locations of the scripts can be configured in
\path{/etc/xen/xend-config.sxp}.

\begin{description} 

\item[network:] This script is called whenever \xend is started or
stopped to respectively initialise or tear down the Xen virtual
network. In the default configuration initialisation creates the
bridge `xen-br0' and moves eth0 onto that bridge, modifying the
routing accordingly. When \xend exits, it deletes the Xen bridge and
removes eth0, restoring the normal IP and routing configuration.

%% In configurations where the bridge already exists, this script could
%% be replaced with a link to \path{/bin/true} (for instance).

\item[vif-bridge:] This script is called for every domain virtual
interface and can configure firewalling rules and add the vif 
to the appropriate bridge. By default, this adds and removes 
VIFs on the default Xen bridge.

\end{description} 

For more complex network setups (e.g. where routing is required or
integrate with existing bridges) these scripts may be replaced with
customised variants for your site's preferred configuration.

%% There are two possible types of privileges:  IO privileges and
%% administration privileges.

\section{Driver Domain Configuration} 

I/O privileges can be assigned to allow a domain to directly access
PCI devices itself.  This is used to support driver domains.

Setting backend privileges is currently only supported in SXP format
config files.  To allow a domain to function as a backend for others,
somewhere within the {\tt vm} element of its configuration file must
be a {\tt backend} element of the form {\tt (backend ({\em type}))}
where {\tt \em type} may be either {\tt netif} or {\tt blkif},
according to the type of virtual device this domain will service.
%% After this domain has been built, \xend will connect all new and
%% existing {\em virtual} devices (of the appropriate type) to that
%% backend.

Note that a block backend cannot currently import virtual block
devices from other domains, and a network backend cannot import
virtual network devices from other domains.  Thus (particularly in the
case of block backends, which cannot import a virtual block device as
their root filesystem), you may need to boot a backend domain from a
ramdisk or a network device.

Access to PCI devices may be configured on a per-device basis.  Xen
will assign the minimal set of hardware privileges to a domain that
are required to control its devices.  This can be configured in either
format of configuration file:

\begin{itemize}
\item SXP Format: Include device elements of the form: \\
\centerline{  {\tt (device (pci (bus {\em x}) (dev {\em y}) (func {\em z})))}} \\
  inside the top-level {\tt vm} element.  Each one specifies the address
  of a device this domain is allowed to access ---
  the numbers {\em x},{\em y} and {\em z} may be in either decimal or
  hexadecimal format.
\item Flat Format: Include a list of PCI device addresses of the
  format: \\ 
\centerline{{\tt pci = ['x,y,z', ...]}} \\ 
where each element in the
  list is a string specifying the components of the PCI device
  address, separated by commas.  The components ({\tt \em x}, {\tt \em
  y} and {\tt \em z}) of the list may be formatted as either decimal
  or hexadecimal.
\end{itemize}

%% \section{Administration Domains}

%% Administration privileges allow a domain to use the `dom0
%% operations' (so called because they are usually available only to
%% domain 0).  A privileged domain can build other domains, set scheduling
%% parameters, etc.

% Support for other administrative domains is not yet available...  perhaps
% we should plumb it in some time





\section{Scheduler Configuration}
\label{s:sched} 


Xen offers a boot time choice between multiple schedulers.  To select
a scheduler, pass the boot parameter {\em sched=sched\_name} to Xen,
substituting the appropriate scheduler name.  Details of the schedulers
and their parameters are included below; future versions of the tools
will provide a higher-level interface to these tools.

It is expected that system administrators configure their system to
use the scheduler most appropriate to their needs.  Currently, the BVT
scheduler is the recommended choice. 

\subsection{Borrowed Virtual Time}

{\tt sched=bvt} (the default) \\ 

BVT provides proportional fair shares of the CPU time.  It has been
observed to penalise domains that block frequently (e.g. I/O intensive
domains), but this can be compensated for by using warping. 

\subsubsection{Global Parameters}

\begin{description}
\item[ctx\_allow]
  the context switch allowance is similar to the `quantum'
  in traditional schedulers.  It is the minimum time that
  a scheduled domain will be allowed to run before being
  pre-empted. 
\end{description}

\subsubsection{Per-domain parameters}

\begin{description}
\item[mcuadv]
  the MCU (Minimum Charging Unit) advance determines the
  proportional share of the CPU that a domain receives.  It
  is set inversely proportionally to a domain's sharing weight.
\item[warp]
  the amount of `virtual time' the domain is allowed to warp
  backwards
\item[warpl]
  the warp limit is the maximum time a domain can run warped for
\item[warpu]
  the unwarp requirement is the minimum time a domain must
  run unwarped for before it can warp again
\end{description}

\subsection{Atropos}

{\tt sched=atropos} \\

Atropos is a soft real time scheduler.  It provides guarantees about
absolute shares of the CPU, with a facility for sharing
slack CPU time on a best-effort basis. It can provide timeliness
guarantees for latency-sensitive domains.

Every domain has an associated period and slice.  The domain should
receive `slice' nanoseconds every `period' nanoseconds.  This allows
the administrator to configure both the absolute share of the CPU a
domain receives and the frequency with which it is scheduled. 

%%  When
%% domains unblock, their period is reduced to the value of the latency
%% hint (the slice is scaled accordingly so that they still get the same
%% proportion of the CPU).  For each subsequent period, the slice and
%% period times are doubled until they reach their original values.

Note: don't overcommit the CPU when using Atropos (i.e. don't reserve
more CPU than is available --- the utilisation should be kept to
slightly less than 100\% in order to ensure predictable behaviour).

\subsubsection{Per-domain parameters}

\begin{description}
\item[period] The regular time interval during which a domain is
  guaranteed to receive its allocation of CPU time.
\item[slice]
  The length of time per period that a domain is guaranteed to run
  for (in the absence of voluntary yielding of the CPU). 
\item[latency]
  The latency hint is used to control how soon after
  waking up a domain it should be scheduled.
\item[xtratime] This is a boolean flag that specifies whether a domain
  should be allowed a share of the system slack time.
\end{description}

\subsection{Round Robin}

{\tt sched=rrobin} \\

The round robin scheduler is included as a simple demonstration of
Xen's internal scheduler API.  It is not intended for production use. 

\subsubsection{Global Parameters}

\begin{description}
\item[rr\_slice]
  The maximum time each domain runs before the next
  scheduling decision is made.
\end{description}












\chapter{Build, Boot and Debug options} 

This chapter describes the build- and boot-time options 
which may be used to tailor your Xen system. 

\section{Xen Build Options}

Xen provides a number of build-time options which should be 
set as environment variables or passed on make's command-line.  

\begin{description} 
\item[verbose=y] Enable debugging messages when Xen detects an unexpected condition.
Also enables console output from all domains.
\item[debug=y] 
Enable debug assertions.  Implies {\bf verbose=y}.
(Primarily useful for tracing bugs in Xen).       
\item[debugger=y] 
Enable the in-Xen debugger. This can be used to debug 
Xen, guest OSes, and applications.
\item[perfc=y] 
Enable performance counters for significant events
within Xen. The counts can be reset or displayed
on Xen's console via console control keys.
\item[trace=y] 
Enable per-cpu trace buffers which log a range of
events within Xen for collection by control
software. 
\end{description} 

\section{Xen Boot Options}
\label{s:xboot}

These options are used to configure Xen's behaviour at runtime.  They
should be appended to Xen's command line, either manually or by
editing \path{grub.conf}.

\begin{description}
\item [ignorebiostables ] 
 Disable parsing of BIOS-supplied tables. This may help with some
 chipsets that aren't fully supported by Xen. If you specify this
 option then ACPI tables are also ignored, and SMP support is
 disabled. 

\item [noreboot ] 
 Don't reboot the machine automatically on errors.  This is
 useful to catch debug output if you aren't catching console messages
 via the serial line. 

\item [nosmp ] 
 Disable SMP support.
 This option is implied by `ignorebiostables'. 

\item [noacpi ] 
 Disable ACPI tables, which confuse Xen on some chipsets.
 This option is implied by `ignorebiostables'. 

\item [watchdog ] 
 Enable NMI watchdog which can report certain failures. 

\item [noht ] 
 Disable Hyperthreading. 

\item [badpage=$<$page number$>$,$<$page number$>$, \ldots ] 
 Specify a list of pages not to be allocated for use 
 because they contain bad bytes. For example, if your
 memory tester says that byte 0x12345678 is bad, you would
 place `badpage=0x12345' on Xen's command line. 

\item [com1=$<$baud$>$,DPS,$<$io\_base$>$,$<$irq$>$
 com2=$<$baud$>$,DPS,$<$io\_base$>$,$<$irq$>$ ] \mbox{}\\ 
 Xen supports up to two 16550-compatible serial ports.
 For example: `com1=9600, 8n1, 0x408, 5' maps COM1 to a
 9600-baud port, 8 data bits, no parity, 1 stop bit,
 I/O port base 0x408, IRQ 5.
 If the I/O base and IRQ are standard (com1:0x3f8,4;
 com2:0x2f8,3) then they need not be specified. 

\item [console=$<$specifier list$>$ ] 
 Specify the destination for Xen console I/O.
 This is a comma-separated list of, for example:
\begin{description}
 \item[vga]  use VGA console and allow keyboard input
 \item[com1] use serial port com1
 \item[com2H] use serial port com2. Transmitted chars will
   have the MSB set. Received chars must have
   MSB set.
 \item[com2L] use serial port com2. Transmitted chars will
   have the MSB cleared. Received chars must
   have MSB cleared.
\end{description}
 The latter two examples allow a single port to be
 shared by two subsystems (e.g. console and
 debugger). Sharing is controlled by MSB of each
 transmitted/received character.
 [NB. Default for this option is `com1,vga'] 

\item [conswitch=$<$switch-char$><$auto-switch-char$>$ ] 
 Specify how to switch serial-console input between
 Xen and DOM0. The required sequence is CTRL-$<$switch-char$>$
 pressed three times. Specifying the backtick character 
 disables switching.
 The $<$auto-switch-char$>$ specifies whether Xen should
 auto-switch input to DOM0 when it boots --- if it is `x'
 then auto-switching is disabled.  Any other value, or
 omitting the character, enables auto-switching.
 [NB. default switch-char is `a'] 

\item [nmi=xxx ] 
 Specify what to do with an NMI parity or I/O error. \\
 `nmi=fatal':  Xen prints a diagnostic and then hangs. \\
 `nmi=dom0':   Inform DOM0 of the NMI. \\
 `nmi=ignore': Ignore the NMI. 

\item [dom0\_mem=xxx ] 
 Set the amount of memory (in kB) to be allocated to domain0.  

\item [tbuf\_size=xxx ] 
 Set the size of the per-cpu trace buffers, in pages
 (default 1).  Note that the trace buffers are only
 enabled in debug builds.  Most users can ignore
 this feature completely. 

\item [sched=xxx ] 
 Select the CPU scheduler Xen should use.  The current
 possibilities are `bvt' (default), `atropos' and `rrobin'. 
 For more information see Section~\ref{s:sched}. 

\item [pci\_dom0\_hide=(xx.xx.x)(yy.yy.y)\ldots ] 
Hide selected PCI devices from domain 0 (for instance, to stop it
taking ownership of them so that they can be driven by another
domain).  Device IDs should be given in hex format.  Bridge devices do
not need to be hidden --- they are hidden implicitly, since guest OSes
do not need to configure them.
\end{description} 



\section{XenLinux Boot Options}

In addition to the standard Linux kernel boot options, we support: 
\begin{description} 
\item[xencons=xxx ] Specify the device node to which the Xen virtual
console driver is attached. The following options are supported:
\begin{center}
\begin{tabular}{l}
`xencons=off': disable virtual console \\ 
`xencons=tty': attach console to /dev/tty1 (tty0 at boot-time) \\
`xencons=ttyS': attach console to /dev/ttyS0
\end{tabular}
\end{center}
The default is ttyS for dom0 and tty for all other domains.
\end{description} 



\section{Debugging}
\label{s:keys} 

Xen has a set of debugging features that can be useful to try and
figure out what's going on. Hit 'h' on the serial line (if you
specified a baud rate on the Xen command line) or ScrollLock-h on the
keyboard to get a list of supported commands.

If you have a crash you'll likely get a crash dump containing an EIP
(PC) which, along with an \path{objdump -d image}, can be useful in
figuring out what's happened.  Debug a Xenlinux image just as you
would any other Linux kernel.

%% We supply a handy debug terminal program which you can find in
%% \path{/usr/local/src/xen-2.0.bk/tools/misc/miniterm/}
%% This should be built and executed on another machine that is connected
%% via a null modem cable. Documentation is included.
%% Alternatively, if the Xen machine is connected to a serial-port server
%% then we supply a dumb TCP terminal client, {\tt xencons}.




\chapter{Further Support}

If you have questions that are not answered by this manual, the
sources of information listed below may be of interest to you.  Note
that bug reports, suggestions and contributions related to the
software (or the documentation) should be sent to the Xen developers'
mailing list (address below).

\section{Other Documentation}

For developers interested in porting operating systems to Xen, the
{\em Xen Interface Manual} is distributed in the \path{docs/}
directory of the Xen source distribution.  

%Various HOWTOs are available in \path{docs/HOWTOS} but this content is
%being integrated into this manual.

\section{Online References}

The official Xen web site is found at:
\begin{quote}
{\tt http://www.cl.cam.ac.uk/netos/xen/}
\end{quote}

This contains links to the latest versions of all on-line 
documentation (including the lateset version of the FAQ). 

\section{Mailing Lists}

There are currently three official Xen mailing lists:

\begin{description}
\item[xen-devel@lists.sourceforge.net] Used for development
discussions and requests for help.  Subscribe at: \\
\path{http://lists.sourceforge.net/mailman/listinfo/xen-devel}
\item[xen-announce@lists.sourceforge.net] Used for announcements only.
Subscribe at: \\
\path{http://lists.sourceforge.net/mailman/listinfo/xen-announce}
\item[xen-changelog@lists.sourceforge.net]  Changelog feed
from the unstable and 2.0 trees - developer oriented.  Subscribe at: \\
\path{http://lists.sourceforge.net/mailman/listinfo/xen-changelog}
\end{description}

Although there is no specific user support list, the developers try to
assist users who post on xen-devel.  As the bulk of traffic on this
list increases, a dedicated user support list may be introduced.

\appendix


\chapter{Installing Xen / XenLinux on Debian}

The Debian project provides a tool called \path{debootstrap} which
allows a base Debian system to be installed into a filesystem without
requiring the host system to have any Debian-specific software (such
as \path{apt}. 

Here's some info how to install Debian 3.1 (Sarge) for an unprivileged
Xen domain:

\begin{enumerate}
\item Set up Xen 2.0 and test that it's working, as described earlier in
      this manual.

\item Create disk images for root-fs and swap (alternatively, you
      might create dedicated partitions, LVM logical volumes, etc. if
      that suits your setup).
\begin{small}\begin{verbatim}  
dd if=/dev/zero of=/path/diskimage bs=1024k count=size_in_mbytes
dd if=/dev/zero of=/path/swapimage bs=1024k count=size_in_mbytes
\end{verbatim}\end{small}
      If you're going to use this filesystem / disk image only as a
      `template' for other vm disk images, something like 300 MB should
      be enough.. (of course it depends what kind of packages you are
      planning to install to the template)

\item Create the filesystem and initialise the swap image
\begin{small}\begin{verbatim}
mkfs.ext3 /path/diskimage
mkswap /path/swapimage
\end{verbatim}\end{small}

\item Mount the disk image for installation
\begin{small}\begin{verbatim}
mount -o loop /path/diskimage /mnt/disk
\end{verbatim}\end{small}

\item Install \path{debootstrap}

Make sure you have debootstrap installed on the host.  If you are
running Debian sarge (3.1 / testing) or unstable you can install it by
running \path{apt-get install debootstrap}.  Otherwise, it can be
downloaded from the Debian project website.

\item Install Debian base to the disk image:
\begin{small}\begin{verbatim}
debootstrap --arch i386 sarge /mnt/disk  \
            http://ftp.<countrycode>.debian.org/debian
\end{verbatim}\end{small}

You can use any other Debian http/ftp mirror you want.

\item When debootstrap completes successfully, modify settings:
\begin{small}\begin{verbatim}
chroot /mnt/disk /bin/bash
\end{verbatim}\end{small}

Edit the following files using vi or nano and make needed changes:
\begin{small}\begin{verbatim}
/etc/hostname
/etc/hosts
/etc/resolv.conf
/etc/network/interfaces
/etc/networks
\end{verbatim}\end{small}

Set up access to the services, edit:
\begin{small}\begin{verbatim}
/etc/hosts.deny
/etc/hosts.allow
/etc/inetd.conf
\end{verbatim}\end{small}

Add Debian mirror to:   
\begin{small}\begin{verbatim}
/etc/apt/sources.list
\end{verbatim}\end{small}

Create fstab like this:
\begin{small}\begin{verbatim}
/dev/sda1       /       ext3    errors=remount-ro       0       1
/dev/sda2       none    swap    sw                      0       0
proc            /proc   proc    defaults                0       0
\end{verbatim}\end{small}

Logout

\item      Unmount the disk image
\begin{small}\begin{verbatim}
umount /mnt/disk
\end{verbatim}\end{small}

\item Create Xen 2.0 configuration file for the new domain. You can
        use the example-configurations coming with Xen as a template.

        Make sure you have the following set up:
\begin{small}\begin{verbatim}
disk = [ 'file:/path/diskimage,sda1,w', 'file:/path/swapimage,sda2,w' ]
root = "/dev/sda1 ro"
\end{verbatim}\end{small}

\item Start the new domain
\begin{small}\begin{verbatim}
xm create -f domain_config_file
\end{verbatim}\end{small}

Check that the new domain is running:
\begin{small}\begin{verbatim}
xm list
\end{verbatim}\end{small}

\item   Attach to the console of the new domain.
        You should see something like this when starting the new domain:

\begin{small}\begin{verbatim}
Started domain testdomain2, console on port 9626
\end{verbatim}\end{small}
        
        There you can see the ID of the console: 26. You can also list
        the consoles with \path{xm consoles} (ID is the last two
        digits of the port number.)

        Attach to the console:

\begin{small}\begin{verbatim}
xm console 26
\end{verbatim}\end{small}

        or by telnetting to the port 9626 of localhost (the xm console
        program works better).

\item   Log in and run base-config

        As a default there's no password for the root.

        Check that everything looks OK, and the system started without
        errors.  Check that the swap is active, and the network settings are
        correct.

        Run \path{/usr/sbin/base-config} to set up the Debian settings.

        Set up the password for root using passwd.

\item     Done. You can exit the console by pressing \path{Ctrl + ]}

\end{enumerate}

If you need to create new domains, you can just copy the contents of
the `template'-image to the new disk images, either by mounting the
template and the new image, and using \path{cp -a} or \path{tar} or by
simply copying the image file.  Once this is done, modify the
image-specific settings (hostname, network settings, etc).

\chapter{Installing Xen / XenLinux on Redhat or Fedora Core}

When using Xen / XenLinux on a standard Linux distribution there are
a couple of things to watch out for:

Note that, because domains>0 don't have any privileged access at all,
certain commands in the default boot sequence will fail e.g. attempts
to update the hwclock, change the console font, update the keytable
map, start apmd (power management), or gpm (mouse cursor).  Either
ignore the errors (they should be harmless), or remove them from the
startup scripts.  Deleting the following links are a good start:
{\path{S24pcmcia}}, {\path{S09isdn}},
{\path{S17keytable}}, {\path{S26apmd}},
{\path{S85gpm}}.

If you want to use a single root file system that works cleanly for
both domain 0 and unprivileged domains, a useful trick is to use
different 'init' run levels. For example, use
run level 3 for domain 0, and run level 4 for other domains. This
enables different startup scripts to be run in depending on the run
level number passed on the kernel command line.

If using NFS root files systems mounted either from an
external server or from domain0 there are a couple of other gotchas.
The default {\path{/etc/sysconfig/iptables}} rules block NFS, so part
way through the boot sequence things will suddenly go dead.

If you're planning on having a separate NFS {\path{/usr}} partition, the
RH9 boot scripts don't make life easy - they attempt to mount NFS file
systems way to late in the boot process. The easiest way I found to do
this was to have a {\path{/linuxrc}} script run ahead of
{\path{/sbin/init}} that mounts {\path{/usr}}:

\begin{quote}
\begin{small}\begin{verbatim}
 #!/bin/bash
 /sbin/ipconfig lo 127.0.0.1
 /sbin/portmap
 /bin/mount /usr
 exec /sbin/init "$@" <>/dev/console 2>&1
\end{verbatim}\end{small}
\end{quote}

%$ XXX SMH: font lock fix :-)  

The one slight complication with the above is that
{\path{/sbin/portmap}} is dynamically linked against
{\path{/usr/lib/libwrap.so.0}} Since this is in
{\path{/usr}}, it won't work. This can be solved by copying the
file (and link) below the /usr mount point, and just let the file be
'covered' when the mount happens.

In some installations, where a shared read-only {\path{/usr}} is
being used, it may be desirable to move other large directories over
into the read-only {\path{/usr}}. For example, you might replace
{\path{/bin}}, {\path{/lib}} and {\path{/sbin}} with
links into {\path{/usr/root/bin}}, {\path{/usr/root/lib}}
and {\path{/usr/root/sbin}} respectively. This creates other
problems for running the {\path{/linuxrc}} script, requiring
bash, portmap, mount, ifconfig, and a handful of other shared
libraries to be copied below the mount point --- a simple
statically-linked C program would solve this problem.




\chapter{Glossary of Terms}

\begin{description}
\item[Atropos]             One of the CPU schedulers provided by Xen.
                           Atropos provides domains with absolute shares
                           of the CPU, with timeliness guarantees and a
                           mechanism for sharing out `slack time'.

\item[BVT]                 The BVT scheduler is used to give proportional
                           fair shares of the CPU to domains.

\item[Exokernel]           A minimal piece of privileged code, similar to
                           a {\bf microkernel} but providing a more
                           `hardware-like' interface to the tasks it
                           manages.  This is similar to a paravirtualising
                           VMM like {\bf Xen} but was designed as a new
                           operating system structure, rather than
                           specifically to run multiple conventional OSs.

\item[Domain]              A domain is the execution context that
                           contains a running {\bf virtual machine}.
                           The relationship between virtual machines
                           and domains on Xen is similar to that between
                           programs and processes in an operating
                           system: a virtual machine is a persistent
                           entity that resides on disk (somewhat like
                           a program).  When it is loaded for execution,
                           it runs in a domain.  Each domain has a
                           {\bf domain ID}.

\item[Domain 0]            The first domain to be started on a Xen
                           machine.  Domain 0 is responsible for managing
                           the system.

\item[Domain ID]           A unique identifier for a {\bf domain},
                           analogous to a process ID in an operating
                           system.

\item[Full virtualisation] An approach to virtualisation which
                           requires no modifications to the hosted
                           operating system, providing the illusion of
                           a complete system of real hardware devices.

\item[Hypervisor]          An alternative term for {\bf VMM}, used
                           because it means `beyond supervisor',
                           since it is responsible for managing multiple
                           `supervisor' kernels.

\item[Live migration]      A technique for moving a running virtual
                           machine to another physical host, without
                           stopping it or the services running on it.

\item[Microkernel]         A small base of code running at the highest
                           hardware privilege level.  A microkernel is
                           responsible for sharing CPU and memory (and
                           sometimes other devices) between less
                           privileged tasks running on the system.
                           This is similar to a VMM, particularly a
                           {\bf paravirtualising} VMM but typically
                           addressing a different problem space and
                           providing different kind of interface.

\item[NetBSD/Xen]          A port of NetBSD to the Xen architecture.

\item[Paravirtualisation]  An approach to virtualisation which requires
                           modifications to the operating system in
                           order to run in a virtual machine.  Xen
                           uses paravirtualisation but preserves
                           binary compatibility for user space
                           applications.

\item[Shadow pagetables]   A technique for hiding the layout of machine
                           memory from a virtual machine's operating
                           system.  Used in some {\bf VMMs} to provide
                           the illusion of contiguous physical memory,
                           in Xen this is used during
                           {\bf live migration}.

\item[Virtual Machine]     The environment in which a hosted operating
                           system runs, providing the abstraction of a
                           dedicated machine.  A virtual machine may
                           be identical to the underlying hardware (as
                           in {\bf full virtualisation}, or it may
                           differ, as in {\bf paravirtualisation}.

\item[VMM]                 Virtual Machine Monitor - the software that
                           allows multiple virtual machines to be
                           multiplexed on a single physical machine.

\item[Xen]                 Xen is a paravirtualising virtual machine
                           monitor, developed primarily by the
                           Systems Research Group at the University
                           of Cambridge Computer Laboratory.

\item[XenLinux]            Official name for the port of the Linux kernel
                           that runs on Xen.

\end{description}


\end{document}


%% Other stuff without a home

%% Instructions Re Python API

%% Other Control Tasks using Python
%% ================================

%% A Python module 'Xc' is installed as part of the tools-install
%% process. This can be imported, and an 'xc object' instantiated, to
%% provide access to privileged command operations:

%% # import Xc
%% # xc = Xc.new()
%% # dir(xc)
%% # help(xc.domain_create)

%% In this way you can see that the class 'xc' contains useful
%% documentation for you to consult.

%% A further package of useful routines (xenctl) is also installed:

%% # import xenctl.utils
%% # help(xenctl.utils)

%% You can use these modules to write your own custom scripts or you can
%% customise the scripts supplied in the Xen distribution.



% Explain about AGP GART


%% If you're not intending to configure the new domain with an IP address
%% on your LAN, then you'll probably want to use NAT. The
%% 'xen_nat_enable' installs a few useful iptables rules into domain0 to
%% enable NAT. [NB: We plan to support RSIP in future]




%% Installing the file systems from the CD
%% =======================================

%% If you haven't got an existing Linux installation onto which you can
%% just drop down the Xen and Xenlinux images, then the file systems on
%% the CD provide a quick way of doing an install. However, you would be
%% better off in the long run doing a proper install of your preferred
%% distro and installing Xen onto that, rather than just doing the hack
%% described below:

%% Choose one or two partitions, depending on whether you want a separate
%% /usr or not. Make file systems on it/them e.g.: 
%%   mkfs -t ext3 /dev/hda3
%%   [or mkfs -t ext2 /dev/hda3 && tune2fs -j /dev/hda3 if using an old
%% version of mkfs]

%% Next, mount the file system(s) e.g.:
%%   mkdir /mnt/root && mount /dev/hda3 /mnt/root
%%   [mkdir /mnt/usr && mount /dev/hda4 /mnt/usr]
  
%% To install the root file system, simply untar /usr/XenDemoCD/root.tar.gz:
%%   cd /mnt/root && tar -zxpf /usr/XenDemoCD/root.tar.gz

%% You'll need to edit /mnt/root/etc/fstab to reflect your file system
%% configuration. Changing the password file (etc/shadow) is probably a
%% good idea too.

%% To install the usr file system, copy the file system from CD on /usr,
%% though leaving out the "XenDemoCD" and "boot" directories:
%%   cd /usr && cp -a X11R6 etc java libexec root src bin dict kerberos local sbin tmp doc include lib man share /mnt/usr

%% If you intend to boot off these file systems (i.e. use them for
%% domain 0), then you probably want to copy the /usr/boot directory on
%% the cd over the top of the current symlink to /boot on your root
%% filesystem (after deleting the current symlink) i.e.:
%%   cd /mnt/root ; rm boot ; cp -a /usr/boot .
